\documentclass[11pt]{article}
\usepackage{../homework}
\profname{Dr. Kurt Hinterbichler}
\classcode{PHYS 365 General Relativity}
\assignnum{Assignment 03}
\duedate{03 November 2026}

\begin{document}
\begin{problem}{1}
Show that
\begin{equation}
  \label{eq:p1-covariant-commutator}
  [\nabla_\mu,\nabla_\nu]V^\alpha
  = R^\alpha{}_{\lambda\mu\nu}V^\lambda - T^\lambda{}_{\mu\nu}\nabla_\lambda V^\alpha,
\end{equation}
where the Riemann curvature and torsion are the following tensors defined in terms of the connection:
\begin{equation}
  \label{eq:p1-torsion}
  T^\lambda{}_{\mu\nu} = \Gamma^\lambda{}_{\mu\nu}-\Gamma^\lambda{}_{\nu\mu},
\end{equation}
\begin{equation}
  \label{eq:p1-riemann}
  R^\alpha{}_{\beta\mu\nu}
  = \partial_\mu\Gamma^\alpha{}_{\nu\beta}-\partial_\nu\Gamma^\alpha{}_{\mu\beta}
    +\Gamma^\alpha{}_{\mu\lambda}\Gamma^\lambda{}_{\nu\beta}
    -\Gamma^\alpha{}_{\nu\lambda}\Gamma^\lambda{}_{\mu\beta}.
\end{equation}
(Recall the expression for the covariant derivative of a general tensor
\begin{equation}
  \label{eq:p1-general-covariant-derivative}
  \begin{aligned}
    \nabla_\alpha T^{\mu\nu\cdots}{}_{\beta\gamma\cdots}
    = {}& \partial_\alpha T^{\mu\nu\cdots}{}_{\beta\gamma\cdots}
      + \Gamma^\mu{}_{\alpha\lambda}T^{\lambda\nu\cdots}{}_{\beta\gamma\cdots}
      + \Gamma^\nu{}_{\alpha\lambda}T^{\mu\lambda\cdots}{}_{\beta\gamma\cdots}+\cdots \\
    &-\Gamma^\lambda{}_{\alpha\beta}T^{\mu\nu\cdots}{}_{\lambda\gamma\cdots}
      -\Gamma^\lambda{}_{\alpha\gamma}T^{\mu\nu\cdots}{}_{\beta\lambda\cdots}-\cdots.)
  \end{aligned}
\end{equation}
\end{problem}

\begin{problem}{2}
From now on we stick to the Levi-Civita connection.
Show that the differential Bianchi identity \(\nabla_{[\rho}R_{\sigma\lambda]\mu\nu}=0\) implies the contracted Bianchi identity
\begin{equation}
  \label{eq:p2-contracted-bianchi}
  \nabla^\nu G_{\mu\nu}=0,
\end{equation}
where \(G_{\mu\nu}=R_{\mu\nu}-\frac12 Rg_{\mu\nu}\) is the Einstein tensor.
\end{problem}

\begin{problem}{3}
Consider Euclidean \(\mathbb{R}^2\) with the standard cartesian coordinates \(x^1,x^2\) and metric \(ds^2=(dx^1)^2+(dx^2)^2\).
Since the metric is constant in these coordinates, the connection vanishes \(\Gamma^\rho{}_{\mu\nu}=0\), and the curvature vanishes, \(R^\rho{}_{\sigma\mu\nu}=0\).
Now change to polar coordinates \(r,\theta\)
\begin{equation}
  \label{eq:p3-polar-coordinates}
  x^1=r\cos\theta,\qquad x^2=r\sin\theta.
\end{equation}
\begin{enumerate}[label=(\alph*)]
  \item Using the transformation rule we derived in class for the connection,
  \begin{equation}
    \label{eq:p3-connection-transformation}
    \Gamma'{}^{\lambda'}{}_{\mu'\nu'}(y)
    =\frac{\partial x^\mu}{\partial y^{\mu'}}
      \frac{\partial x^\nu}{\partial y^{\nu'}}
      \left.\frac{\partial y^{\lambda'}}{\partial x^\lambda}\right|_{x(y)}
      \Gamma^\lambda{}_{\mu\nu}(x(y))
      +\frac{\partial^2 x^\alpha}{\partial y^{\mu'}\partial y^{\nu'}}
      \left.\frac{\partial y^{\lambda'}}{\partial x^\alpha}\right|_{x(y)},
  \end{equation}
  convert the cartesian connection into polar coordinates.
  \item Check your answer to part a. by computing the connection directly from the metric in polar coordinates,
  \begin{equation}
    \label{eq:p3-polar-metric}
    ds^2=dr^2+r^2d\theta^2,
  \end{equation}
  using the formula for the connection in terms of the metric:
  \begin{equation}
    \label{eq:p3-levi-civita-connection}
    \Gamma^\rho{}_{\mu\nu}
    =\frac12 g^{\rho\lambda}(\partial_\mu g_{\nu\lambda}+\partial_\nu g_{\mu\lambda}-\partial_\lambda g_{\mu\nu}).
  \end{equation}
  \item Compute the Riemann curvature in polar coordinates directly from the connection using the expression \cref{eq:p1-riemann} and verify that it vanishes (as it must since this is just flat space).
\end{enumerate}
\end{problem}

\begin{problem}{4}
Recall the parallel transport equation which tells you how to parallel transport a vector \(V^\mu\) along a curve \(x^\mu(\lambda)\),
\begin{equation}
  \label{eq:p4-parallel-transport}
  \frac{d}{d\lambda}V^\mu+\Gamma^\mu{}_{\nu\rho}\frac{dx^\nu}{d\lambda}V^\rho=0.
\end{equation}
Consider a curve in Euclidean \(\mathbb{R}^2\) given in cartesian coordinates by \(x^1(\lambda)=\cos\lambda\), \(x^2(\lambda)=\sin\lambda\), \(\lambda\in[0,\frac\pi2]\).
This curve traces out a quarter of the unit circle.
Consider a vector which starts out as a unit vector pointing in the \(x^2\) direction: \(\left.V^\mu\right|_{\lambda=0}=(0,1)\).
\begin{enumerate}[label=(\alph*)]
  \item Parallel transport on flat space conforms to our intuition that a vector remains constant as it is transported along the curve, so we expect \(V^\mu(\lambda)=(0,1)\) for all \(\lambda\) along the curve.
  Verify that this indeed satisfies the parallel transport equation \cref{eq:p4-parallel-transport}.
  \item Consider now the same setup in polar coordinates.
  The equation for the curve is now
  \begin{equation}
    \label{eq:p4-polar-curve}
    r(\lambda)=1,\qquad\theta(\lambda)=\lambda,
  \end{equation}
  and the initial vector is pointed in the \(\theta\) direction: \(\left.V^\mu\right|_{\lambda=0}=(0,1)\).
  Solve the parallel transport equation \cref{eq:p4-parallel-transport} directly in polar coordinates with these initial conditions, using the connection coefficients you found in problem 1.
  Show that the resulting vector points in the \(r\) direction at the end of the curve, \(\left.V^\mu\right|_{\lambda=\pi/2}=(1,0)\), and is in fact the same vector you get by parallel transporting in cartesian coordinates.
\end{enumerate}
\end{problem}

\begin{problem}{5}
Consider the covariant laplacian of a scalar function \(\phi\),
\begin{equation}
  \label{eq:p5-covariant-laplacian}
  \nabla^2\phi\equiv\nabla_\mu\nabla^\mu\phi=g^{\mu\nu}\nabla_\mu\nabla_\nu\phi.
\end{equation}
In cartesian coordinates on \(\mathbb{R}^2\), this reduces to the ordinary Laplacian \(\frac{\partial^2\phi}{\partial(x^1)^2}+\frac{\partial^2\phi}{\partial(x^2)^2}\).
Using the connection coefficients you computed in problem 3b from the metric \cref{eq:p3-polar-metric}, find the Laplacian in polar coordinates and compare with the well-known expression
\begin{equation}
  \label{eq:p5-polar-laplacian}
  \frac1r\frac{d}{dr}\left(r\frac{d\phi}{dr}\right)+\frac1{r^2}\frac{d^2\phi}{d\theta^2}.
\end{equation}
\end{problem}

\begin{problem}{6. Derivation of the Schwarzschild solution}
The Schwarzschild solution was the first non-trivial exact solution to the Einstein field equations, derived by Karl Schwarzschild while he was in the German army fighting on the Russian front during World War I.
Here you will go through his derivation.

The most general static time-independent metric can be written in the form
\begin{equation}
  \label{eq:p6-static-metric}
  ds^2=-A(r)dt^2+B(r)dr^2+r^2d\Omega^2,
\end{equation}
where \(d\Omega^2=d\theta^2+\sin^2\theta\,d\phi^2\) is the metric on a 2-sphere.
When \(B(r)=A(r)=1\), this is flat Minkowski space in spherical coordinates.
We would like our solution to approach flat space at infinity, so we demand the boundary conditions
\begin{equation}
  \label{eq:p6-boundary-conditions}
  A(r)\xrightarrow{r\to\infty}1,\qquad B(r)\xrightarrow{r\to\infty}1.
\end{equation}
We are interested in vacuum solutions, i.e. solutions that exists outside any matter, so we set \(T_{\mu\nu}=0\) and look for solutions to the vacuum Einstein equations
\begin{equation}
  \label{eq:p6-vacuum-einstein}
  G_{\mu\nu}=0.
\end{equation}
\begin{enumerate}[label=(\alph*)]
  \item Show that the vacuum Einstein equations \cref{eq:p6-vacuum-einstein} are equivalent to the vanishing of the Ricci tensor:
  \begin{equation}
    \label{eq:p6-vacuum-ricci}
    R_{\mu\nu}=0.
  \end{equation}
  \item Compute the non-vanishing components of the Ricci tensor.
  You should find
  \begin{equation}
    \label{eq:p6-ricci-components}
    \begin{aligned}
      R_{tt}&=\frac{A''}{2B}-\frac{A'B'}{4B^2}-\frac{A'^2}{4AB}+\frac{A'}{rB},\\
      R_{rr}&=-\frac{A''}{2A}+\frac{A'B'}{4AB}+\frac{A'^2}{4A^2}+\frac{B'}{rB},\\
      R_{\theta\theta}&=-\frac{rA'}{2AB}+\frac{rB'}{2B^2}-\frac1B+1,\\
      R_{\phi\phi}&=\sin^2\theta\,R_{\theta\theta}.
    \end{aligned}
  \end{equation}
  \item Form the combination \(\frac BA R_{tt}+R_{rr}=0\) and show that it reduces to the equation
  \begin{equation}
    \label{eq:p6-log-derivative-relation}
    \frac{A'}A+\frac{B'}B=0.
  \end{equation}
  Show that this implies \(B=c/A\) for some constant \(c\).
  Based on the boundary conditions \cref{eq:p6-boundary-conditions}, argue that we must have \(c=1\) so that
  \begin{equation}
    \label{eq:p6-reciprocal-relation}
    B=\frac1A.
  \end{equation}
  \item Plug \cref{eq:p6-reciprocal-relation} into \(R_{\theta\theta}\) and show that this gives the equation
  \begin{equation}
    \label{eq:p6-radial-equation}
    -rA'-A+1=0.
  \end{equation}
  Solve this to obtain
  \begin{equation}
    \label{eq:p6-radial-solution}
    A(r)=1-\frac{R_S}r,
  \end{equation}
  with \(R_S\) an integration constant.
  \item Show that \cref{eq:p6-radial-solution} and \cref{eq:p6-reciprocal-relation} together solve \(R_{tt}=0\) and \(R_{rr}=0\) and thus give a solution to the full Einstein equations.
  \item We now have the metric
  \begin{equation}
    \label{eq:p6-schwarzschild-parameter-metric}
    ds^2=-\left(1-\frac{R_S}r\right)dt^2+\frac{1}{1-\frac{R_S}r}dr^2+r^2d\Omega^2,
  \end{equation}
  which is a solution for any \(R_S\).
  To fix the meaning of \(R_S\), expand the metric around flat space \(g_{\mu\nu}=\eta_{\mu\nu}+h_{\mu\nu}\).
  For small \(h_{\mu\nu}\), we saw in class that we can identify the Newtonian potential as \(h_{00}=-2\Phi\).
  Given that the Newtonian potential for a spherical mass is \(\Phi=-\frac{GM}r\), expand \cref{eq:p6-schwarzschild-parameter-metric} near \(r\to\infty\) where \(h_{\mu\nu}\) is small and the Newtonian approximation is good, and use this to identify
  \begin{equation}
    \label{eq:p6-schwarzschild-radius}
    R_S=2GM.
  \end{equation}
  This is the Schwarzschild radius.
\end{enumerate}
The final metric for the Schwarzschild solution thus reads
\begin{equation}
  \label{eq:p6-schwarzschild-metric}
  ds^2=-\left(1-\frac{2GM}r\right)dt^2+\frac{1}{1-\frac{2GM}r}dr^2+r^2d\Omega^2.
\end{equation}
\end{problem}
\end{document}
